\section{Conclusion}
\label{sec:concl}

We have proposed a new non-perturbative method to renormalize lattice
composite operators. It can be used in all cases and is particularly
useful when
the Ward identity method is not applicable. This method avoids the need to
perform any calculations using lattice perturbation theory in the
computation of physical quantities from lattice simulations.
 The success
of our proposal is subject to the existence, at current values of
$\beta$, of a window between the non-perturbative region  at low momenta
 and the region of large momenta, where discretization errors become
important. This  limitation is however, common to all methods
(with the exception of the Ward identity method, which can only be  applied
to a few cases, corresponding to finite operators). The window, necessary
to implement  the renormalization programme described here, seems
 to exists already
at $\beta=6.0$ and we expect that the range of momenta, useful
for non-perturbative renormalization, will become larger  as $\beta$
increases. We are planning to apply the approach presented in this paper to the
renormalization of the $\Delta S=2$ four fermion operator (\ref{eq:ds2b}) and
of
the heavy-light axial vector current in the static theory.
